\section{Limitations}
\textbf{Compute.} Every arm was trained for \Iters{} iterations, chosen by the budget rule on a free Kaggle T4 (\SPerIter\,s per iteration, \AmpText). The FreeNeRF paper trains far longer on a different base (mip-NeRF in JAX); absolute PSNR values are therefore lower than published ones and our 100-view reference is not converged. $\rho$ is defined against a reference trained with the same budget, which makes it a fair within-study ratio but not a comparison with published dense-view numbers.
\textbf{Statistics.} The main comparison uses 3 seeds and 25 test views per scene; the ablation and the view sweep use one seed, so their differences are reported without seed intervals and should be read as indicative.
\textbf{Scope.} Two synthetic objects with white backgrounds; one regulariser and one comparator; default hyper-parameters from the paper, not tuned for nerfstudio.

\section{Contribution statement and AI declaration}
\textbf{Contributions.} \emph{Waris (1024170155) and Tanishq (1024170156)} designed the project: the problem framing and proposal, the detailed module-by-module mapping of what to build, and the explanation flowcharts of the pipeline (data split, method, evaluation, hypotheses and controls) that the implementation follows. They chose the regulariser, the evaluation protocol and the self-captured scene, captured the photos, and Tanishq built and ran an independent plain-PyTorch implementation of the frequency mask used as a cross-check (Fig.~\ref{fig:crosscheck}). They reviewed the code and the results, and take responsibility for every claim made here.
\textbf{AI declaration.} We used Claude (Anthropic) for the writing work: it wrote the code and tests from our design, ran and monitored the Kaggle training jobs on our account, and drafted the text of this report and of the slides from the logged results, which we reviewed. Every number in this report is produced by \texttt{run\_all.sh} from the run logs committed to the public repository; no number was typed by hand.
